% Source: https://artofproblemsolving.com/community/c3h2727251_cool_property
% Stable identity: aops-topic-2727251-q02
The equation becomes $2abc=99a+9b$, i.e. $a(2bc-99)=9b$. We need $2bc>99$, so $bc\ge50$, and then $a=\dfrac{9b}{2bc-99}$ must be an integer from $1$ to $9$. The digit pairs with $bc\ge50$ are $(6,9)$, $(7,8)$, $(7,9)$, $(8,7)$, $(8,8)$, $(8,9)$ and $(9,c)$ for $c=6,7,8,9$. Checking each:
\begin{itemize}
\item $b=9$: $a=\dfrac{81}{18c-99}$, which is $9$ for $c=6$ and $3$ for $c=7$; $c=8,9$ give $\dfrac{81}{45}$ and $\dfrac{81}{63}$.
\item $b=8$: $a=\dfrac{72}{16c-99}$ is $\dfrac{72}{13},\dfrac{72}{29},\dfrac{72}{45}$ for $c=7,8,9$: never an integer.
\item $b=7$: $a=\dfrac{63}{14c-99}$ is $\dfrac{63}{13},\dfrac{63}{27}$ for $c=8,9$: never an integer.
\item $b=6$, $c=9$: $a=\dfrac{54}{9}=6$.
\end{itemize}
The numbers are $\boxed{397,\ 669,\ 996}$ (e.g. $3+9+7+2\cdot189=397$).
